\chapter{Le support à l'infrastructure : deux semaines de « fast »}

\section{Le principe du fast}

En parallèle du projet d'alerting, j'ai assuré pendant deux semaines le rôle de « fast » au sein de
l'équipe Infrastructure. Le fast désigne la personne désignée, sur une période donnée, pour répondre
aux demandes et aux problèmes remontés par les équipes de développement. Pour l'infrastructure,
ces demandes sont variées : provisionner une nouvelle ressource (un cluster Kafka, une base de
données), résoudre un problème de connexion réseau, débloquer un déploiement ou accompagner
une équipe sur une configuration. Ce rôle est exigeant car il demande de la réactivité, une bonne
connaissance transverse de la plateforme et la capacité à passer rapidement d'un contexte technique
à un autre. Il constitue aussi une formidable occasion d'apprentissage, en confrontant à la diversité
réelle des besoins des équipes.

\section{Exemples de demandes traitées}

\subsection{Ajout d'une ressource AWS : un Valkey pour l'équipe PKD}

L'une des demandes traitées consistait à provisionner une instance Valkey — un moteur de cache et
de stockage clé-valeur compatible Redis, issu d'un fork open source — pour l'équipe PKD. J'ai réalisé
cet ajout en Infrastructure as Code avec Terraform : définition de la ressource, paramétrage réseau
(sous-réseaux, groupes de sécurité), puis soumission via Pull Request pour revue et application par la
CI/CD. Cette tâche m'a permis de me familiariser avec le cycle de provisionnement d'un datastore au
sein de l'organisation AWS de Winamax.

\subsection{Résolution de problèmes de connexion}

J'ai également traité plusieurs problèmes de connexion remontés par les équipes. Ce type d'incident,
fréquent dans une infrastructure cloud segmentée en de nombreux comptes et VPC, nécessite une
démarche méthodique : identifier les extrémités concernées, vérifier le routage réseau, les règles de
filtrage, la résolution DNS et la configuration des VPN. Ces interventions m'ont appris à diagnostiquer
un problème de bout en bout entre services répartis sur plusieurs environnements.

\subsection{Migration d'un master Puppet vers AWS}

La tâche la plus structurante de cette période a été la migration d'un master Puppet. Trois serveurs
pointaient encore vers un ancien master hébergé hors d'AWS ; je les ai reconfigurés pour qu'ils
pointent vers le master hébergé sur AWS, afin de centraliser la gestion de configuration. Cette
migration a impliqué de reconfigurer les agents, de valider la nouvelle communication avec le master
et de s'assurer de l'absence de régression. Cette tâche illustre bien la valeur du rôle de fast :
contribuer à la dette technique et à la simplification de l'infrastructure.

\section{Apports de cette expérience}

Ces deux semaines de fast ont été particulièrement formatrices. Elles m'ont permis de développer une
vision transverse de l'infrastructure de Winamax, de gagner en autonomie dans le diagnostic et la
résolution d'incidents, et de renforcer mes compétences sur des outils clés comme Terraform, AWS
et Puppet. Elles ont aussi mis en évidence l'importance de la communication : bien comprendre le
besoin d'une équipe, tenir informé le demandeur de l'avancement et documenter la résolution sont
des réflexes essentiels. Enfin, cette immersion dans le quotidien opérationnel de l'équipe a donné un
sens concret au projet d'alerting : c'est en étant soi-même confronté aux sollicitations des équipes
que l'on mesure toute la valeur d'un système de notification pertinent et bien calibré.
